\subsection{A monotone sequence converging to the full constant}
The analytic cubic sequence below has limit $1/\mathcal N(b)$.
The reduction also permits an unconditional sequence converging to the
full Grothendieck constant, without assuming that a cubic operator is
optimal.

\begin{proposition}\label{prop:universal-monotone}
Let $\mathcal A_n$ consist of all nonzero real matrices of dimensions
at most $n$ whose entries lie in
$2^{-n}\mathbb Z\cap[-1,1]$. Choose a matrix $S_n\in\mathcal A_n$
maximizing $G$, with a fixed lexicographic rule for ties. Clear its
dyadic denominator and apply \eqref{eq:clone} and \eqref{eq:balanced},
obtaining $U_n$. Then every $U_n$ belongs to the simultaneous reduced
class and
\[
 G(U_n)\le G(U_{n+1}),\qquad
 \lim_{n\to\infty}G(U_n)=G_{\mathbb R}.
\]
\end{proposition}
\begin{proof}
Each candidate set is finite and the sets are nested. Both norms are
positive on a nonzero matrix, and the exact reductions preserve their
ratio, so the maxima exist and increase monotonically.
For any fixed nonzero real matrix $A$, homogeneity permits the
normalization $\max_{ij}|A_{ij}|\le1$. Its dyadic approximants $A^{(n)}$
belong to $\mathcal A_n$ for sufficiently large $n$ and have total
entry error tending to zero. The two norms satisfy
\[
 |C(A)-C(B)|\le\sum_{ij}|A_{ij}-B_{ij}|,
 \qquad |V(A)-V(B)|\le\sum_{ij}|A_{ij}-B_{ij}|.
\]
Consequently $G(A^{(n)})\to G(A)$ and
$\liminf_nG(U_n)\ge G(A)$. Taking the supremum over $A$ proves the
asserted limit; the reverse inequality follows from the definition of
$G_{\mathbb R}$.
\end{proof}

This is a finite exact specification at every stage. The scalar norm is
computed by sign enumeration. The vector norm is the maximum of the
rational linear functional
$\sum_{ij}A_{ij}X_{i,m+j}$ on the compact set of positive semidefinite
$(m+n)$-by-$(m+n)$ matrices with diagonal one. Its polynomial
description makes exact algebraic optimization and comparison possible.
Enumeration is computationally expensive. The analytic subfamily gives
controlled candidates and a quantitative limit, while retaining a best
previous candidate prevents a weaker trial from replacing it.
